% ====================================================================
\chapter{Pruebas y validación}
\label{cap:pruebas_validacion}
% ====================================================================

En el presente capítulo se expone la estrategia de verificación y validación del sistema de traducción en tiempo pseudo-real de LSA. Se detallan el plan de pruebas adoptado, la ejecución de los experimentos sobre los distintos componentes y el análisis cuantitativo y cualitativo de los resultados obtenidos frente a los requerimientos del proyecto. El objetivo central consiste en evidenciar la aplicación de criterios técnicos para evaluar el funcionamiento, el desempeño y la calidad del software, con una distinción clara entre planificación, ejecución y análisis de resultados.

% ====================================================================
\section{Plan de pruebas}
\label{sec:plan_pruebas}
% ====================================================================

El plan de pruebas establece la estrategia sistemática para comprobar el correcto funcionamiento, el desempeño en tiempo real y la calidad del software desarrollado.

\subsection{Alcance y tipos de pruebas}
La evaluación del sistema abarca desde la recolección de datos y el entrenamiento de modelos de aprendizaje profundo hasta la integración completa de los componentes en una videollamada web. Se aplicaron los siguientes tipos de pruebas:
\begin{itemize}
    \item \textbf{Pruebas unitarias y de componentes:} Verificación del extractor de puntos clave (\textit{landmarks}), la lógica de filtrado gestual (\texttt{RepeatGate}), los clasificadores de señas y la cuantización de los modelos de lenguaje.
    \item \textbf{Pruebas de integración:} Evaluación de las interfaces entre la extensión de navegador (Manifest V3), el motor de inferencia local, las llamadas a API y el cliente de videoconferencia.
    \item \textbf{Pruebas de rendimiento y latencia:} Medición de tiempos de respuesta por etapa del canal de procesamiento, consumo de memoria en cliente y servidor, y tasa de fotogramas por segundo.
    \item \textbf{Pruebas de sistema de extremo a extremo (\textit{End-to-End}):} Validación del circuito completo de traducción (LSA a español y español a LSA) en entornos de llamada reales.
    \item \textbf{Pruebas de robustez y casos límite:} Evaluación del clasificador frente a cambios de iluminación, distorsiones ópticas de cámara web, presencia de múltiples personas en el fondo y variabilidad en la velocidad de articulación gestual.
\end{itemize}

\subsection{Entornos de evaluación}
Las pruebas se ejecutaron en tres entornos representativos:
\begin{enumerate}
    \item \textbf{Estación de desarrollo (GPU):} Utilizada para la evaluación estática fuera de línea, el entrenamiento de modelos y las pruebas de velocidad máxima de inferencia.
    \item \textbf{Equipo objetivo de cliente (Laptop sin GPU):} Configuración de gama media estándar orientada a verificar el comportamiento en tiempo real, el consumo de CPU/RAM y la latencia perceptible por el usuario final. Sus componentes están en la Sección~\ref{sec:laptop_prueba}.
    \item \textbf{Infraestructura en la nube:} Instancia de Google Cloud Run (configuración de 4 vCPU y 8\,GiB de memoria RAM) conectada mediante HTTPS para los servicios remotos de traducción semántica y reproducción de avatar.
\end{enumerate}

\subsection{Equipo portátil de las pruebas de cliente}
\label{sec:laptop_prueba}
La sesión de la Sección~\ref{sec:perfilado_cliente_recursos} corrió en una laptop con Windows 11 Home. El Cuadro~\ref{tab:laptop_prueba} reúne los componentes de ese equipo. El nombre comercial del equipo no quedó registrado. La toma del 30 de octubre de 2026, con RAM y CPU del navegador, corrió en la estación de escritorio y está en la Sección~\ref{sec:hw_comparacion}.

\begin{table}[H]
    \centering
    \small
    \caption{Laptop de las pruebas de cliente.}
    \label{tab:laptop_prueba}
    \begin{tabular}{@{}ll@{}}
        \toprule
        \textbf{Componente} & \textbf{Valor} \\
        \midrule
        Procesador & Intel Core i7-10510U a 1{,}80\,GHz \\
        Memoria RAM & 8\,GB a 2667\,MT/s \\
        Gráficos & Intel UHD Graphics, 128\,MB \\
        Almacenamiento & 477\,GB \\
        Sistema operativo & Windows 11 Home \\
        \bottomrule
    \end{tabular}
\end{table}

\subsection{Criterios de aceptación y trazabilidad}
Cada caso de prueba se vinculó de manera directa con los requerimientos funcionales y no funcionales del sistema. En el Cuadro~\ref{tab:trazabilidad_pruebas} se sintetizan los criterios de aceptación principales definidos para la validación.

% TODO: CHEQUEAR FRANN

\begin{table}[H]
    \centering
    \small
    \caption{Matriz de trazabilidad entre requerimientos y criterios de aceptación para pruebas.}
    \label{tab:trazabilidad_pruebas}
    \begin{tabularx}{\textwidth}{l X X}
        \toprule
        \textbf{Módulo} & \textbf{Requerimiento asociado} & \textbf{Criterio de aceptación} \\
        \midrule
        Clasificador Gestual & RF-01, RF-02 & Exactitud Top-1 en cámara real $> 75\%$ y Top-3 $> 90\%$. Latencia de inferencia $< 10$\,ms. \\
        Traductor Semántico & RF-03, RF-06 & Exactitud sintáctica/semántica $> 85\%$, BLEU $> 90$, BERTScore $> 98$. \\
        Motor Local & RF-02, RF-09 & Supresión de falsos positivos y rebotes sin descartar repeticiones intencionales. \\
        Integración Meet & RF-08 & Sin interrupción de la pista de video principal; renderizado de subtítulos sin pérdidas. \\
        \bottomrule
    \end{tabularx}
\end{table}

% ====================================================================
\section{Resultados de la verificación}
\label{sec:resultados_verificacion}
% ====================================================================

Esta sección presenta la evidencia empírica obtenida a partir de la ejecución del plan de pruebas sobre cada componente del sistema.

\subsection{Evaluación del clasificador gestual (TinySkeletonClassifier)}
\label{sec:res_clasificador}

El modelo de clasificación gestual se evaluó tanto en entornos estáticos fuera de línea como en pruebas en vivo frente a cámara web en la estación local.

\subsubsection{Evolución del corpus y vocabulario}
La incorporación de clases al modelo avanzó de manera paulatina conforme se consolidaron las muestras del dataset propio de LSA. En el Cuadro~\ref{tab:evolucion_corpus} se detalla la progresión del vocabulario y las muestras recopiladas.

\begin{table}[H]
    \centering
    \small
    \caption{Progresión cronológica del conjunto de datos y vocabulario de LSA.}
    \label{tab:evolucion_corpus}
    \begin{tabular}{lcccc}
        \toprule
        \textbf{Nro} & \textbf{Clases} & \textbf{Muestras/Clase} & \textbf{Muestras Entrenamiento} & \textbf{Muestras Validación} \\
        \midrule
        1 & 91 & 50 & 3\,640 & 910 \\
        2 & 94 & 60 & 4\,512 & 1\,128 \\
        3 & 97 & 60 & 4\,656 & 1\,164 \\
        4 & 100 & 60 & 4\,800 & 1\,200 \\
        5 & 100 & 100 & 8\,000 & 2\,000 \\
        \bottomrule
    \end{tabular}
\end{table}

\subsubsection{Comparativa entre evaluación fuera de línea y cámara en vivo}
En el Cuadro~\ref{tab:brecha_offline_camara} se expone la diferencia observada entre la exactitud de validación sobre el dataset estático y el rendimiento real obtenido frente a la cámara web.

\begin{table}[H]
    \centering
    \small
    \caption{Exactitud de validación fuera de línea frente al acierto en cámara. Las filas de 91 a 97 clases corresponden a una toma por seña. Las de 100 clases corresponden a diez tomas por seña.}
    \label{tab:brecha_offline_camara}
    \begin{tabularx}{\linewidth}{Xcccc}
        \toprule
        \textbf{Checkpoint} & \textbf{Dataset} & \textbf{Val. Top-1} & \textbf{Cámara Top-1} & \textbf{Cámara Top-3} \\
        \midrule
        16 fotogramas, 91 clases & 1 & 97{,}36\,\% & 61{,}5\,\% & 83{,}5\,\% \\
        32 fotogramas, 128/4/2 & 1 & 96{,}48\,\% & 71{,}4\,\% & 87{,}9\,\% \\
        32 fotogramas, 64/2/1 & 1 & 97{,}36\,\% & 67{,}0\,\% & 90{,}1\,\% \\
        Optuna v1, 32 fotogramas & 1 & 97{,}69\,\% & 66{,}0\,\% & 84{,}6\,\% \\
        Optuna v2, 12 fotogramas, 94 clases & 2 & 97{,}78\,\% & 80{,}9\,\% & 96{,}8\,\% \\
        Baseline 16f, luz ambiental, 94 clases & 2 & 98{,}40\,\% & 67{,}0\,\% & 89{,}4\,\% \\
        Baseline 16f, iluminación dirigida, 94 clases & 2 & 98{,}40\,\% & 86{,}2\,\% & 95{,}7\,\% \\
        Baseline 16f, 97 clases & 3 & 98{,}28\,\% & 77{,}3\,\% & 94{,}8\,\% \\
        100 clases, 10 tomas, luz adecuada & 4 & 99{,}00\,\% & 99{,}6\,\% & 100{,}0\,\% \\
        100 clases, 10 tomas, luz reducida & 5 & 99{,}00\,\% & 94{,}6\,\% & 98{,}7\,\% \\
        \bottomrule
    \end{tabularx}
\end{table}

En las iteraciones de 94 clases, el control de la iluminación del entorno elevó el acierto en cámara del 67{,}0\,\% al 86{,}2\,\% en Top-1 sobre la misma arquitectura, sin modificar los pesos del modelo.

\subsubsection{Resolución temporal: 16 frente a 32 fotogramas}
El Cuadro~\ref{tab:frames_julio} sintetiza los experimentos realizados para determinar la longitud óptima de la ventana temporal de entrada.

\begin{table}[H]
    \centering
    \small
    \caption{Evaluación de la resolución temporal en julio de 2026 (91 clases, 50 videos por clase).}
    \label{tab:frames_julio}
    \begin{tabularx}{\textwidth}{l X ccc}
        \toprule
        \textbf{Iteración} & \textbf{Arquitectura} & \textbf{Val. Top-1} & \textbf{Cám. Top-1} & \textbf{Cám. Top-3} \\
        \midrule
        A, 16 frames & 128 dim / 4 cabezas / 2 capas & 97{,}36\,\% & 61{,}5\,\% & 83{,}5\,\% \\
        B, 32 frames & 128 dim / 4 cabezas / 2 capas & 96{,}48\,\% & 71{,}4\,\% & 87{,}9\,\% \\
        C, 32 frames & 64 dim / 2 cabezas / 1 capa   & 97{,}36\,\% & 67{,}0\,\% & 90{,}1\,\% \\
        \bottomrule
    \end{tabularx}
\end{table}

A pesar de las ganancias iniciales observadas con 32 fotogramas, la incorporación de aumentaciones avanzadas de datos en agosto confirmó que la ventana de 16 fotogramas ofreció mayor capacidad de generalización sobre señas estáticas y redujo la redundancia temporal.

\subsubsection{Búsqueda de hiperparámetros con Optuna}

La exploración del espacio de hiperparámetros se estructuró en dos estudios compuestos por 40 ensayos cada uno, evaluados mediante validación cruzada de tres pliegues con poda de ensayos no prometedores. La función objetivo priorizó la minimización de la pérdida media de validación fuera de línea, sin incluir la exactitud directa frente a la cámara web.

El primer estudio abarcó un vocabulario de 91 clases con una ventana fija de 32 fotogramas. La configuración seleccionada por la búsqueda conservó la dimensión de 128 y 2 capas, al tiempo que duplicó las cabezas de atención a 8, redujo la tasa de \textit{dropout} a $0{,}30$, fijó el tamaño de lote en 16 e introdujo una tasa de aprendizaje de $3{,}77\times 10^{-5}$. Tras el reentrenamiento final, este modelo obtuvo una exactitud de validación del 97{,}69\,\% y una pérdida de $0{,}087$. En las pruebas frente a cámara, las métricas registraron un 66{,}0\,\% en Top-1 y un 84{,}6\,\% en Top-3. Dichos valores se situaron 5,4 puntos porcentuales por debajo de la iteración manual B en Top-1 y 5,5 puntos por debajo de la iteración C en Top-3.

El segundo estudio amplió el catálogo a 94 clases e incorporó la resolución temporal (ventanas de 8, 12 y 16 fotogramas) dentro del espacio de búsqueda. El ensayo de mayor rendimiento seleccionó una ventana de 12 fotogramas, 256 dimensiones, 2 cabezas de atención, 2 capas y un \textit{dropout} de 0{,}57. En la evaluación fuera de línea, la pérdida descendió a 0{,}080 y la exactitud alcanzó el 97{,}78\,\%, lo que resultó 0,62 puntos porcentuales inferior al \textit{baseline} de 16 fotogramas (98{,}40\,\%). Bajo condiciones controladas de iluminación en cámara, la exactitud Top-1 se registró en 80{,}9\,\%, frente al 86{,}2\,\% obtenido por el \textit{baseline}. Esta degradación afectó de forma directa al alfabeto dactilológico, donde la exactitud en las letras cayó 14,8 puntos porcentuales (con fallos específicos en \texttt{A}, \texttt{N}, \texttt{V} e \texttt{Y}). Si bien la métrica Top-3 mostró un leve incremento en sesiones independientes, una evaluación apareada sobre tomas idénticas descartó la existencia de una ventaja estadística significativa.

\begin{table}[H]
    \centering
    \small
    \caption{Comparativa entre los ensayos de Optuna y las arquitecturas de referencia. Las cifras en cámara del segundo estudio corresponden a idénticas condiciones de iluminación que el \textit{baseline} de 16 fotogramas.}
    \label{tab:optuna_vs_baseline}
    \begin{tabularx}{\textwidth}{Xlccc}
        \toprule
        \textbf{Modelo} & \textbf{Arquitectura} & \textbf{Val. Top-1} & \textbf{Cám. Top-1} & \textbf{Cám. Top-3} \\
        \midrule
        Optuna v1, 32 fotogramas & 128 / 8 / 2 & 97{,}69\,\% & 66{,}0\,\% & 84{,}6\,\% \\
        Iteración B, 32 fotogramas & 128 / 4 / 2 & 96{,}48\,\% & 71{,}4\,\% & 87{,}9\,\% \\
        Optuna v2, 12 fotogramas & 256 / 2 / 2 & 97{,}78\,\% & 80{,}9\,\% & 96{,}8\,\% \\
        \textit{Baseline} 16 frames, misma luz & 128 / 4 / 2 & 98{,}40\,\% & 86{,}2\,\% & 95{,}7\,\% \\
        \bottomrule
    \end{tabularx}
\end{table}

El descarte de las configuraciones producidas por Optuna obedeció a tres razones fundamentales:

\begin{enumerate}
    \item \textbf{Disparidad de dominio:} La pérdida de validación sobre secuencias previamente curadas no predice el comportamiento real frente a la cámara, puesto que las fluctuaciones de iluminación, los errores de encuadre y la inestabilidad en la extracción de puntos articulares permanecen ajenos al proceso de validación cruzada.
    \item \textbf{Aislamiento de variables en la captura:} Las comparaciones iniciales solaparon los cambios arquitectónicos con variaciones en las condiciones de luz. Una vez estabilizada la iluminación ambiental, se constató que el incremento en el rendimiento provino de la calidad de la captura y no de los hiperparámetros.
    \item \textbf{Criterio de desempeño del producto:} El diseño del sistema prioriza la exactitud Top-1 en vivo, con especial énfasis en el deletreo dactilológico. La arquitectura base (16 fotogramas, 128 dimensiones, 4 cabezas de atención y 2 capas) superó a los dos modelos generados por la optimización automática y constituyó la solución integrada en el prototipo final.
\end{enumerate}


\subsubsection{Impacto de la corrección de relación de aspecto}
La estandarización del marco de captura desde proporciones $4:3$ hacia $16:9$ con bandas de relleno generó un impacto positivo directo en la clasificación, según se detalla en el Cuadro~\ref{tab:aspect}.

\begin{table}[H]
    \centering
    \small
    \caption{Efecto de la corrección de relación de aspecto sobre 36 tomas de letras y dígitos.}
    \label{tab:aspect}
    \begin{tabular}{lc}
        \toprule
        \textbf{Configuración de entrada} & \textbf{Exactitud Top-1} \\
        \midrule
        Captura original en relación 4:3 & 66{,}7\,\% \\
        Corrección analítica en eje $y$ ($\times 1{,}333$) & 88{,}9\,\% \\
        Corrección simulada en ejes $y$ y $z$ & 75{,}0\,\% \\
        \bottomrule
    \end{tabular}
\end{table}

La transformación del eje $y$ aportó un incremento de 22,2 puntos porcentuales en exactitud. Este ajuste corrigió confusiones sistemáticas, como el caso de la letra \texttt{I} clasificada previamente como \texttt{ojo}.

\subsubsection{Evaluación en cámara del modelo de 100 clases}
\label{sec:eval_100_camara}
\label{sec:protocolo_camara}

El prototipo integrado carga el modelo \texttt{model\_new\_2026\_10\_20}: 100 clases, 16 fotogramas, 128 dimensiones, 4 cabezas y 2 capas. El entrenamiento usó 100 clips por clase (8\,000 secuencias de entrenamiento y 2\,000 de validación) y alcanzó una exactitud de validación del 99{,}0\,\%, con pérdida de $0{,}028$. Ese corte es por clip dentro de cada clase, como se describe en la Sección~\ref{sec:tinyskeleton}. Las mismas tres personas pueden estar en el entrenamiento y en la validación.

La prueba en cámara consistió en dos sesiones sobre ese mismo archivo. La mano, el encuadre y el extractor se mantuvieron fijos; solo varió la iluminación de la habitación. Quien ejecuta el script anota la etiqueta de luz adecuada o reducida, dado que el programa no mide la iluminancia. Cada sesión recorrió las 100 señas con diez tomas consecutivas (1\,000 inferencias por condición), en modo automático y con mano derecha, sin omitir tomas.

La cámara se solicita a $1280\times 720$ y entrega $640\times 480$. El script rellena bandas hasta $853\times 480$, el mismo marco 16:9 de los videos de entrenamiento. El extractor corresponde a MediaPipe Holistic en navegador (\texttt{@mediapipe/holistic@0.5.1675471629}): pose y dos manos, sin malla facial, vector de 225 coordenadas y resampleo a 16 fotogramas. Entre inferencias se estableció un segundo de espera.

El Cuadro~\ref{tab:eval_100_luz} resume los resultados. Las dos condiciones de iluminación cumplen el criterio de aceptación del plan (Top-1 mayor al 75\,\% y Top-3 mayor al 90\,\%).

\begin{table}[H]
    \centering
    \small
    \caption{Evaluación en cámara del modelo (\texttt{model\_new\_2026\_10\_20}). Diez tomas por seña, misma persona y mismo encuadre.}
    \label{tab:eval_100_luz}
    \begin{tabular}{lcc}
        \toprule
        \textbf{Métrica} & \textbf{Luz adecuada} & \textbf{Luz reducida} \\
        \midrule
        Top-1 & 99{,}6\,\% (996/1000) & 94{,}6\,\% (946/1000) \\
        Top-3 & 100{,}0\,\% (1000/1000) & 98{,}7\,\% (987/1000) \\
        Señas con 10 aciertos sobre 10 & 97 & 77 \\
        Errores que restan & 4 & 54 \\
        Confianza media en aciertos & $0{,}979$ & $0{,}967$ \\
        Confianza media en errores & $0{,}632$ & $0{,}771$ \\
        \bottomrule
    \end{tabular}
\end{table}

Con luz adecuada, 97 señas registran diez aciertos en diez tomas. Restan cuatro errores y en todos ellos la seña solicitada ocupa el segundo lugar: \texttt{A} se confunde dos veces con \texttt{6}; \texttt{ahora\_hoy} una vez con \texttt{tuyo}; \texttt{repetir} una vez con \texttt{esposo\_a}. El subconjunto de 37 letras, eñe y dígitos alcanza un 99{,}5\,\% de Top-1.

Con luz reducida el Top-1 desciende 5,0 puntos y 20 señas dejan de completar diez aciertos sobre diez. El Top-3 permanece por encima del 98\,\%. Las caídas más marcadas en señas con diez o nueve aciertos con luz adecuada corresponden a \texttt{T} (3/10, hacia \texttt{I}), \texttt{como} (3/10, hacia \texttt{numero}), \texttt{años} (5/10, hacia \texttt{G}), \texttt{hola} (6/10, hacia \texttt{chau}) y \texttt{ahora\_hoy} (6/10, con destinos diversos). En trece tomas la seña solicitada no entra al Top-3. En luz reducida la confianza media de los errores asciende a $0{,}771$, indicador de que la red se equivoca con mayor certeza que con luz adecuada.

\subsection{Resultados de los módulos de traducción semántica}
\label{sec:res_traduccion}
\label{sec:eval_semantica_cuant}

La evaluación de los modelos de lenguaje reducidos (SLM) cuantizados en formato GGUF (\texttt{Q4\_K\_M}) se ejecutó sobre la partición final del corpus (119 oraciones para el canal directo LSA $\rightarrow$ Español y 102 oraciones para el canal inverso Español $\rightarrow$ LSA).

\subsubsection{Evaluación del canal directo (LSA a Español)}
En el Cuadro~\ref{tab:sem_directo} se presentan las métricas cualitativas y cuantitativas para el canal directo.

\begin{table}[H]
\centering
\footnotesize
\setlength{\tabcolsep}{3.5pt}
\caption{Canal directo (LSA $\rightarrow$ Español), resultados finales (119 oraciones de evaluación). Índices en escala 0--100.}
\label{tab:sem_directo}
\begin{tabular}{@{}lrrrrrrr@{}}
\toprule
\textbf{Modelo} & \textbf{Exacto} & \textbf{BLEU} & \textbf{ROUGE-L} & \textbf{METEOR} & \textbf{BERTScore} & \textbf{COMET} & \textbf{ms} \\
\midrule
Qwen2.5 0,5B & 92,44 & 95,58 & 95,78 & 95,51 & 99,16 & 97,47 & 141,05 \\
Qwen2.5 1,5B & 94,96 & 96,32 & 98,02 & 97,35 & 99,20 & 97,36 & 256,30 \\
Qwen2.5 3B   & 95,80 & 97,88 & 98,39 & 97,84 & 99,20 & 97,56 & 514,38 \\
Llama 3.2 1B & 89,92 & 92,64 & 95,94 & 96,13 & 98,58 & 96,83 & 231,28 \\
SmolLM2 1,7B & 90,76 & 93,53 & 96,48 & 95,97 & 98,85 & 97,12 & 438,31 \\
\bottomrule
\end{tabular}
\end{table}

\subsubsection{Evaluación del canal inverso (Español a LSA)}
En el Cuadro~\ref{tab:sem_inverso} se detallan las métricas correspondientes al módulo de estructuración en glosas para el canal inverso.

\begin{table}[H]
\centering
\footnotesize
\setlength{\tabcolsep}{3.5pt}
\caption{Canal inverso (Español $\rightarrow$ LSA), resultados finales (102 oraciones de evaluación). Índices en escala 0--100.}
\label{tab:sem_inverso}
\begin{tabular}{@{}lrrrrrrr@{}}
\toprule
\textbf{Modelo} & \textbf{Exacto} & \textbf{BLEU} & \textbf{ROUGE-L} & \textbf{METEOR} & \textbf{BERTScore} & \textbf{COMET} & \textbf{ms} \\
\midrule
Qwen2.5 0,5B & 87,25 & 91,29 & 96,20 & 91,25 & 98,86 & 94,28 & 155,97 \\
Qwen2.5 1,5B & 91,18 & 94,14 & 97,17 & 92,75 & 99,22 & 95,98 & 290,82 \\
Qwen2.5 3B   & 89,22 & 92,51 & 96,85 & 92,71 & 99,21 & 95,45 & 537,58 \\
Llama 3.2 1B & 89,22 & 92,66 & 96,36 & 92,89 & 99,13 & 95,29 & 245,70 \\
SmolLM2 1,7B & 77,45 & 82,90 & 91,54 & 86,31 & 97,68 & 92,37 & 458,39 \\
\bottomrule
\end{tabular}
\end{table}

El modelo Qwen2.5 1,5B alcanzó el equilibrio óptimo en ambas direcciones, con un acierto exacto del 94,96\,\% en canal directo y del 91,18\,\% en canal inverso. Esos cuadros y la elección del modelo usan el mismo 15\,\% reservado, descrito en la Sección~\ref{sec:eval_lora}. Las medias de 256,30\,ms y 290,82\,ms salen de la evaluación local con la GPU de la estación de desarrollo. El equipo, el programa y la longitud de cada pasada están en la Sección~\ref{sec:hw_comparacion}.

\subsection{Pruebas de la lógica de control del motor local}
La batería de pruebas diseñada para validar la lógica del filtro \texttt{RepeatGate} y los tiempos de corte de enunciado produjo los resultados consignados en el Cuadro~\ref{tab:bateria_logica}.

\begin{table}[H]
\centering
\small
\caption{Resultados de las pruebas funcionales sobre la lógica de control del motor local.}
\label{tab:bateria_logica}
\begin{tabularx}{\textwidth}{l l X}
\toprule
\textbf{Caso de prueba} & \textbf{Entrada ingresada} & \textbf{Resultado emitido} \\
\midrule
Dígitos repetidos & \texttt{1 1 1 2 2} & Aceptación total: \texttt{1 1 1 2 2} \\
Letra repetida tres veces & \texttt{A A A B} & Filtrado a máximo de dos: \texttt{A A B} \\
Rebote del clasificador & \texttt{HOLA HOLA} & Descarte de duplicado inmediato: \texttt{HOLA} \\
Repetición no consecutiva & \texttt{HOLA VER HOLA} & Conservación de la secuencia completa \\
Glosa repetida descartada a los 3\,s & Glosa repetida & El cierre se desplaza a los 7\,s \\
Movimiento continuo hasta 5\,s & Sin glosa nueva & Cierre de enunciado a los 9\,s \\
Deletreo dactilológico & \texttt{J U A N} & Conversión directa a «Juan.» sin consulta a nube \\
Deletreo con ambigüedad & \texttt{T I N A} & Derivación al traductor por contener \texttt{T} e \texttt{I} \\
\bottomrule
\end{tabularx}
\end{table}

\subsection{Tratamiento de fallas de integración en Google Meet}
Durante la evaluación sobre la interfaz de usuario en llamadas activas se identificaron inconvenientes técnicos que se solucionaron mediante intervenciones específicas de arquitectura, según se detalla en el Cuadro~\ref{tab:fallas_integracion}.

\begin{table}[H]
\centering
\small
\caption{Anomalías detectadas en la integración con Google Meet y soluciones aplicadas.}
\label{tab:fallas_integracion}
\begin{tabularx}{\textwidth}{@{}>{\raggedright\arraybackslash}p{0.42\textwidth} >{\raggedright\arraybackslash}X@{}}
\toprule
\textbf{Falla detectada} & \textbf{Intervención de ingeniería aplicada} \\
\midrule
Bloqueo del inicio de video en Meet por intercepción permanente de cámara. & Activación del reemplazo de la pista únicamente con el motor ILSA encendido. \\
\addlinespace[2pt]
Pérdida de enlace al recargar la extensión con la pestaña de la reunión abierta. & Recarga automática de la pestaña de Meet ante el reinicio del contexto de la extensión. \\
\addlinespace[2pt]
Detección de falsas manos provocada por movimiento del torso. & Exigencia de un número mínimo de cuadros consecutivos con mano detectada antes de abrir el gesto. \\
\addlinespace[2pt]
Interferencia en la captura por movimiento de personas en el fondo. & Delimitación del área de actividad basada exclusivamente en los puntos articulares de las manos. \\
\addlinespace[2pt]
Sobrecarga de CPU en la transferencia de imágenes hacia la ventana \textit{offscreen}. & Reducción de la resolución de los cuadros transferidos a $240$ píxeles de ancho con compresión ajustada. \\
\bottomrule
\end{tabularx}
\end{table}


% ====================================================================
\section{Métricas de calidad}
\label{sec:metricas_calidad}
% ====================================================================

La evaluación objetiva del sistema se consolidó mediante métricas cuantificables agrupadas en cuatro dimensiones principales: desempeño en tiempo real, eficiencia en el uso de recursos, calidad semántica y confiabilidad.

\subsection{Pruebas de latencia y end-to-end en llamada real}
\label{sec:eval_telemetria_vivo}
\label{sec:latencias_pipeline}

Se ejecutó una sesión de pruebas end-to-end sobre la extensión integrada en videollamada web activa, en la laptop de la Sección~\ref{sec:laptop_prueba}. Durante dicha sesión se capturaron eventos en ambos canales de traducción (directo e inverso), con registro de la latencia por etapa, de \texttt{RAM\_mb} y \texttt{CPU\_percent} del proceso ILSA, y de la tasa de fotogramas por segundo (FPS). Ese registro no trae la RAM ni la CPU del navegador.

En el Cuadro~\ref{tab:telemetria_resumen} se sintetizan los valores promedios, mínimos y máximos recolectados durante las ejecuciones de prueba.

\begin{table}[H]
\centering
\small
\caption{Resumen de métricas de perfilado y latencia registradas.}
\label{tab:telemetria_resumen}
\begin{tabularx}{\textwidth}{l X X X}
\toprule
\textbf{Métrica de ejecución} & \textbf{Canal Directo (LSA $\rightarrow$ Esp)} & \textbf{Canal Inverso (Esp $\rightarrow$ LSA)} & \textbf{Unidad} \\
\midrule
Eventos registrados & 16 & 11 & oraciones \\
Pausa de cierre (configurada) & $4126{,}5 \pm 68{,}0$ & $0{,}0$ & ms \\
Tiempo de inferencia remota & $1527{,}9$ ($0{,}0$ -- $2525{,}5$) & $1407{,}3$ ($0{,}0$ -- $2224{,}1$) & ms \\
Tiempo de transporte de red & $299{,}1$ ($0{,}0$ -- $350{,}2$) & $278{,}5$ ($0{,}0$ -- $376{,}1$) & ms \\
Tiempo de entrega / renderizado & $0{,}4$ ($0{,}2$ -- $1{,}0$) & $988{,}9$ ($175{,}3$ -- $2945{,}1$) & ms \\
Latencia total de procesamiento & $1827{,}5$ ($0{,}3$ -- $2853{,}2$) & $2674{,}7$ ($175{,}3$ -- $4707{,}7$) & ms \\
RAM de ILSA & $114{,}4$ ($82{,}3$ -- $222{,}1$) & $247{,}7$ ($200{,}3$ -- $286{,}7$) & MB \\
Tasa de procesamiento (Holistic) & $0{,}1$ -- $12{,}4$ & N/A & FPS \\
Tasa de renderizado cliente & $10{,}8$ -- $40{,}9$ & $32{,}9$ -- $40{,}5$ & FPS \\
Cumplimiento RNF-01 ($< 2{,}0$\,s) & 93{,}8\,\% (15/16) & 90{,}9\,\% (10/11) & aciertos \\
\bottomrule
\end{tabularx}
\end{table}

\subsection{Evaluación del uso de recursos de hardware en el dispositivo cliente}
\label{sec:perfilado_cliente_recursos}

A partir de los datos recopilados en cliente, se analizó el impacto del motor de inferencia local sobre los recursos de cómputo de la laptop objetivo sin GPU dedicada. \texttt{RAM\_mb} y \texttt{CPU\_percent} pertenecen al proceso de Python de ILSA (\texttt{os.getpid()}). \texttt{RAM\_mb} es el working set de ese proceso. \texttt{CPU\_percent} es el porcentaje de CPU del mismo proceso y supera el 100\,\% cuando el proceso usa varios núcleos. El navegador queda fuera de esas dos columnas. Esta toma no registró \texttt{Browser\_RAM\_mb} ni \texttt{Browser\_CPU\_percent}. La extensión puede enviar, aparte, \texttt{js\_heap\_mb}: el heap de JavaScript de la página. Ese heap no es la RAM del proceso del navegador.

En el Cuadro~\ref{tab:perfilado_recursos_log} se comparan los requerimientos de memoria y CPU observados durante la ejecución continua de la extensión.

\begin{table}[H]
\centering
\small
\caption{RAM y CPU de ILSA en la laptop durante la prueba activa. El navegador no entra en este cuadro.}
\label{tab:perfilado_recursos_log}
\setlength{\tabcolsep}{4pt}
\begin{tabular}{@{}lcccc@{}}
\toprule
\textbf{Canal} & \textbf{RAM media (MB)} & \textbf{RAM máx. (MB)} & \textbf{CPU media} & \textbf{CPU máx.} \\
\midrule
Directo & 114{,}36 & 222{,}10 & 9{,}33\,\% & 60{,}10\,\% \\
Inverso & 247{,}67 & 286{,}70 & 40{,}02\,\% & 98{,}40\,\% \\
\bottomrule
\end{tabular}
\end{table}

El análisis de perfilado muestra los siguientes comportamientos:
\begin{itemize}
\item \textbf{Eficiencia en memoria:} La RAM de ILSA en el canal directo descendió desde 222{,}1 MB hasta un régimen de 82{,}3 MB. En el canal inverso se estabilizó cerca de 247{,}7 MB.
\item \textbf{Carga del procesador:} Durante el procesamiento del canal directo la CPU de ILSA permaneció en 9{,}33\,\% de promedio. En el canal inverso, Vosk y el armado de las secuencias llevaron la media a 40{,}02\,\%, con un pico de 98{,}40\,\%.
\end{itemize}

\subsection{Consumo de memoria y huella del sistema en servidor}
\label{sec:memoria_inferencia}
La cuantización en formato GGUF de 4 bits permitió reducir la huella de memoria del modelo de lenguaje a valores compatibles con despliegues eficientes. En el Cuadro~\ref{tab:pesos_cuant} y el Cuadro~\ref{tab:ram_inferencia} se sintetizan el tamaño de los archivos y el consumo estimado de memoria RAM en inferencia con un contexto de 2048 tokens.

\begin{table}[H]
\centering
\footnotesize
\caption{Dimensiones de los archivos cuantizados y compresión alcanzada.}
\label{tab:pesos_cuant}
\begin{tabular}{lrrrr}
\toprule
\textbf{Modelo} & \textbf{16 bits (GiB)} & \textbf{GGUF (MiB)} & \textbf{GGUF (GiB)} & \textbf{bits/parámetro} \\
\midrule
Qwen2.5 0,5B & 0,913 & 379,4 & 0,370 & 6,5 \\
Qwen2.5 1,5B & 2,869 & 940,4 & 0,918 & 5,1 \\
Qwen2.5 3B   & 5,756 & 1,840,5 & 1,797 & 5,0 \\
Llama 3.2 1B & 2,291 & 770,3 & 0,752 & 5,3 \\
SmolLM2 1,7B & 3,185 & 1,006,7 & 0,983 & 4,9 \\
\bottomrule
\end{tabular}
\end{table}


\begin{table}[H]
\centering
\footnotesize
\caption{Estimación del consumo de memoria RAM en inferencia (contexto de 2048 tokens).}
\label{tab:ram_inferencia}
\begin{tabular}{lrrr}
\toprule
\textbf{Modelo} & \textbf{Caché $KV$ (GiB)} & \textbf{16 bits + caché (GiB)} & \textbf{GGUF + caché (GiB)} \\
\midrule
Qwen2.5 0,5B & 0,023 & 0,936 & 0,394 \\
Qwen2.5 1,5B & 0,055 & 2,923 & 0,973 \\
Qwen2.5 3B   & 0,070 & 5,826 & 1,867 \\
Llama 3.2 1B & 0,062 & 2,354 & 0,815 \\
SmolLM2 1,7B & 0,375 & 3,560 & 1,358 \\
\bottomrule
\end{tabular}
\end{table}

La columna de caché $KV$ sale de un contexto de 2048 tokens en fp16, dos bytes por elemento \cite{pope2023efficiently}. El primer 2 separa la clave del valor:

\begin{equation}
KV = 2 , n_{\mathrm{layers}} , n_{\mathrm{kv}} , d_{\mathrm{head}} , 2048 \times 2.
\label{eq:kv_cache}
\end{equation}

$n_{\mathrm{layers}}$ es el número de capas, $n_{\mathrm{kv}}$ el de cabezas de clave y valor, y $d_{\mathrm{head}}$ el largo de cada vector. El Cuadro~\ref{tab:kv_calculo} aplica esa cuenta. El cociente en GiB es el número de bytes dividido por $1024^{3}$. El Cuadro~\ref{tab:ram_inferencia} guarda ese cociente con tres decimales.

\begin{table}[H]
\centering
\footnotesize
\setlength{\tabcolsep}{4pt}
\caption{Datos de la caché $KV$ para 2048 tokens en fp16.}
\label{tab:kv_calculo}
\begin{tabular}{lrrrrr}
\toprule
\textbf{Modelo} & \textbf{Capas} & \textbf{Cabezas $KV$} & $d_{\text{head}}$ & \textbf{Bytes} & \textbf{GiB} \\
\midrule
Qwen2.5 0,5B & 24 & 2 & 64 & 25,165,824 & 0,023 \\
Qwen2.5 1,5B & 28 & 2 & 128 & 58,720,256 & 0,055 \\
Qwen2.5 3B & 36 & 2 & 128 & 75,497,472 & 0,070 \\
Llama 3.2 1B & 16 & 8 & 64 & 67,108,864 & 0,062 \\
SmolLM2 1,7B & 24 & 32 & 64 & 402,653,184 & 0,375 \\
\bottomrule
\end{tabular}
\end{table}

Qwen2.5 0,5B da $0{,}0234375$,GiB y el cuadro lo deja en $0{,}023$. Qwen2.5 1,5B da $0{,}0546875$,GiB y el cuadro lo deja en $0{,}055$. Qwen2.5 3B da $0{,}0703125$,GiB y el cuadro lo deja en $0{,}070$. Llama 3.2 1B da $0{,}0625$,GiB y el cuadro lo deja en $0{,}062$. SmolLM2 1,7B da $0{,}375$,GiB exactos. En Qwen2.5 0,5B la cabeza mide 64. En Qwen2.5 1,5B y 3B mide 128. En Llama 3.2 1B, $2048 / 32 = 64$. En SmolLM2 1,7B, $2048 / 32 = 64$, y las 32 cabezas de atención son también las de $KV$.En la réplica de producción desplegada en Cloud Run, ambos canales cargan Qwen2.5 1,5B con ese contexto de 2048. La columna GGUF más caché da $0{,}973$,GiB para ese modelo, y los dos canales suman $1{,}946$,GiB. Esa suma es el archivo medido más la caché $KV$ en fp16. El buffer de activaciones de la Ecuación~\ref{eq:ram_runtime} queda fuera, y la cifra no es el RSS del proceso. Llama 3.2 1B queda en disco, a $0{,}815$,GiB con su caché por canal, y entra en memoria solo si el modelo principal no carga.

\subsection{Comparativa de rendimiento según arquitectura de hardware}
\label{sec:hw_comparacion}
\label{sec:hardware_objetivo}

La toma del 30 de octubre de 2026 corrió en la estación de desarrollo. Esa PC reúne un procesador AMD Ryzen 7 8700F de ocho núcleos a 4,10\,GHz, 16\,GB de memoria RAM a 4800\,MT/s, una GPU NVIDIA GeForce RTX 5070 de 12\,GB y un almacenamiento de 1,37\,TB. El navegador de esa toma es Microsoft Edge. \texttt{RAM\_mb} y \texttt{CPU\_percent} son el proceso de ILSA. \texttt{Browser\_RAM\_mb} y \texttt{Browser\_CPU\_percent} suman los procesos \texttt{msedge.exe}: 31 en el canal directo, y 30 o 31 en el inverso. El heap de JavaScript de la página es otro número, y el del documento donde corre Holistic es otro.

El Cuadro~\ref{tab:telemetria_pc} usa el tramo del registro que trae esas columnas del navegador. Hay 25 oraciones en el canal directo (20 pasan por el modelo, 3 quedan en forma literal y 2 en forma léxica) y 18 en el inverso (15 por el modelo y 3 léxicas). La cola de inferencia queda en 0\,ms en las dos. La pausa del directo mide de 4000 a 4050\,ms, con media 4025\,ms.

\begin{table}[H]
\centering
\footnotesize
\setlength{\tabcolsep}{3pt}
\caption[Toma del 30 de octubre de 2026 en la estación de escritorio]{Toma del 30 de octubre de 2026 en la estación de escritorio. ILSA es el proceso de Python. El navegador es la suma de los procesos de Edge.}
\label{tab:telemetria_pc}
\begin{tabularx}{\textwidth}{@{}>{\raggedright\arraybackslash}p{4.8cm} X X@{}}
\toprule
\textbf{Métrica} & \textbf{Canal directo} & \textbf{Canal inverso} \\
\midrule
Inferencia del modelo, media (mín.--máx.) & 1521{,}3\,ms (942{,}0 -- 2178{,}1) & 1545{,}9\,ms (1004{,}4 -- 2088{,}8) \\
Inferencia del modelo $\le 2$\,s & 17/20 & 14/15 \\
Transporte del modelo, media (mín.--máx.) & 291{,}5\,ms (267{,}0 -- 361{,}1) & 301{,}1\,ms (272{,}2 -- 403{,}1) \\
Entrega, media (mín.--máx.) & 0{,}1\,ms (0{,}0 -- 0{,}3) & 854{,}3\,ms (61{,}4 -- 1532{,}6) \\
Latencia sin la pausa, media (mín.--máx.) & 1450{,}3\,ms (0{,}1 -- 2445{,}3) & 2393{,}4\,ms (61{,}4 -- 3759{,}9) \\
Espera hasta el subtítulo, media (mín.--máx.) & 5475{,}4\,ms (4007{,}2 -- 6487{,}3) & 2393{,}4\,ms (61{,}4 -- 3759{,}9) \\
RAM de ILSA, media (mín.--máx.) & 551{,}1\,MB (536{,}8 -- 554{,}8) & 664{,}7\,MB (549{,}0 -- 686{,}9) \\
CPU de ILSA, mediana (máximo) & 2{,}4\,\% (746{,}4\,\%) & 3{,}5\,\% (83{,}1\,\%) \\
RAM de Edge, media (mín.--máx.) & 3685{,}0\,MB (3525{,}0 -- 3953{,}5) & 3987{,}2\,MB (3787{,}1 -- 4082{,}7) \\
CPU de Edge, media (mín.--máx.) & 187{,}7\,\% (48{,}8 -- 390{,}8) & 126{,}1\,\% (77{,}0 -- 222{,}3) \\
Heap JS de la página, media (mín.--máx.) & 95{,}4\,MB (76{,}1 -- 124{,}1) & 125{,}8\,MB (84{,}3 -- 174{,}4) \\
Heap JS de Holistic, media (mín.--máx.) & 4{,}8\,MB (1{,}1 -- 8{,}7) & N/A \\
FPS del lienzo & 179{,}7 -- 201{,}3 & 192{,}0 -- 200{,}1 \\
FPS de Holistic & 22{,}4 -- 28{,}9; una muestra en 0 & N/A \\
Clasificador, mediana (máximo) & 1{,}4\,ms (80{,}4\,ms) & N/A \\
Transcripción local, media (mín.--máx.) & N/A & 417{,}4\,ms (261{,}6 -- 1040{,}2) \\
Armado del clip, media (mín.--máx.) & N/A & 37{,}2\,ms (4{,}1 -- 116{,}5) \\
\bottomrule
\end{tabularx}
\end{table}

La meta de 2\,s cubre la inferencia del modelo. En esas filas quedan 17 de 20 y 14 de 15. Las 5 oraciones locales del directo y las 3 del inverso tienen inferencia 0, y el registro las cuenta dentro del techo: 22 de 25 y 17 de 18. De 98 señas del directo, 83 entran a la frase y 15 quedan fuera, con confianza bajo 0,75. En 10 de 18 oraciones del inverso la descarga de poses queda en 0\,ms, con el archivo ya en caché. Las otras 8 tardan de 320{,}9 a 1254{,}4\,ms. Cuatro oraciones tienen una glosa sin archivo y una restaura la Ñ. De 14 locuciones, tres no devuelven texto.

El Cuadro~\ref{tab:perfilado_recursos_log} mide solo ILSA en la laptop. En la PC, ILSA queda cerca de 551\,MB en el directo y de 665\,MB en el inverso. Edge queda cerca de 3,7\,GB y de 4,0\,GB. La CPU de ILSA en la PC tiene mediana 2{,}4\,\% en el directo, con un muestreo de 746{,}4\,\%, y mediana 3{,}5\,\% en el inverso, con máximo 83{,}1\,\%. La CPU de Edge tiene media 187{,}7\,\% y 126{,}1\,\%. El lienzo de la PC va de 179{,}7 a 201{,}3\,FPS en el directo y de 192{,}0 a 200{,}1\,FPS en el inverso. Holistic en la PC, canal directo, va de 22{,}4 a 28{,}9\,FPS. En la laptop el mismo extractor va de 0{,}1 a 12{,}4\,FPS. Esas tasas describen el dibujo y el extractor.
